\section*{Capítulo \ref{ch:g11:plate-motions} --- Las placas y sus movimientos}

\begin{solution}{exo:g11:plate-motions:1}
Una \omterm{def:g11:plate-motions:rigid-plate}{placa rígida} es un trozo de \omterm{def:g11:lithosphere-asthenosphere:lithosphere}{litosfera} que se mueve como un solo bloque:
las distancias entre sus puntos no cambian, y su deformación se concentra en
zonas estrechas a lo largo de sus límites. Los terremotos lo muestran: casi
todos se producen a lo largo de franjas estrechas, los límites de placas, y
muy pocos en el interior de las placas.
\end{solution}

\begin{solution}{exo:g11:plate-motions:2}
\qty{72}{mm} al año son \qty{72}{km} por millón de años. En 30 años la
estación se desplaza $72 \times 30 = \qty{2160}{mm}$, unos \qty{2.2}{m}; en un
millón de años, \qty{72}{km}.
\end{solution}

\begin{solution}{exo:g11:plate-motions:3}
$v_E = 30/10 = \qty{3}{cm}$ al año y $v_N = 40/10 = \qty{4}{cm}$ al año. La
velocidad es $\sqrt{3^2 + 4^2} = \qty{5}{cm}$ al año. El ángulo $\beta$ con el
norte cumple $\tan\beta = 3/4$, luego $\beta \approx \ang{37}$: la estación se
mueve hacia el nordeste, en dirección norte desviada \ang{37} hacia el este.
\end{solution}

\begin{solution}{exo:g11:plate-motions:4}
Una \omterm{def:g11:plate-motions:hotspot-track}{traza de punto caliente} es la línea de volcanes y montes submarinos que
una placa se lleva desde un punto caliente al moverse sobre él, cada vez más
antiguos a medida que se alejan del volcán activo. El volcán más joven,
Kilauea, está en el extremo sudeste de la cadena, en la isla de Hawái; las
edades aumentan hacia el oeste-noroeste. La placa se lleva sus volcanes lejos
del punto caliente en la dirección de su movimiento: la placa del Pacífico se
mueve hacia el oeste-noroeste.
\end{solution}

\begin{solution}{exo:g11:plate-motions:5}
$519/5.1 \approx \qty{102}{km}$ por millón de años, es decir, \qty{102}{mm}
al año, unos \qty{10}{cm} al año.
\end{solution}

\begin{solution}{exo:g11:plate-motions:6}
$v_E = 63/31 \approx \qty{20.3}{mm}$ al año y $v_N = 48/31 \approx
\qty{15.5}{mm}$ al año. La velocidad es $\sqrt{20.3^2 + 15.5^2} \approx
\qty{25.5}{mm}$ al año; $\tan\alpha = 15.5/20.3 \approx 0.76$ da $\alpha
\approx \ang{37}$: la estación se mueve en dirección este desviada \ang{37}
hacia el norte, hacia el nordeste. El Mauna Kea se mueve $72/25.5 \approx 2.8$
veces más deprisa.
\end{solution}

\begin{solution}{exo:g11:plate-motions:7}
(1) En milímetros al año,
\[
  \vec v_{\mathrm{NA/EU}} = (-10.5 - 7.7,\ 20.6 - 16.3) = (-18.2, +4.3).
\] Su módulo es $\sqrt{18.2^2 + 4.3^2} \approx
\qty{18.7}{mm}$ al año, y $\tan\alpha = 4.3/18.2 \approx 0.24$ da una
dirección oeste desviada \ang{13} hacia el norte. (2) Las dos placas se
separan a ambos lados de la dorsal a razón de unos \qty{1.9}{cm} al año: es la
\omterm{def:g11:drift-to-spreading:spreading-rate}{velocidad de expansión total} de la dorsal en Islandia. (3) Concuerda con los
\qty{2}{cm} al año que se citan para el Atlántico Norte.
\end{solution}

\begin{solution}{exo:g11:plate-motions:8}
(1) $v = \pi \times 0.651 \times 6371/180 \approx \qty{72.4}{km}$ por millón
de años, unos \qty{7.2}{cm} al año. (2) $72.4 \times \sin\ang{30} \approx
\qty{36.2}{km}$ por millón de años, unos \qty{3.6}{cm} al año. (3) En el polo
de rotación, si está en la placa. La velocidad cambia de un punto a otro,
desde cero en el polo hasta su valor máximo a \ang{90} de él: una placa entera
no tiene una sola velocidad, solo una velocidad en cada punto (o una media
sobre su superficie).
\end{solution}

\begin{solution}{exo:g11:plate-motions:9}
(1) En milímetros al año,
\[
  \vec v_{\mathrm{IN/EU}} = (40.6 - 24.0,\ 38.6 + 3.3) = (+16.6, +41.9).
\] Su módulo es $\sqrt{16.6^2 + 41.9^2} \approx
\qty{45}{mm}$ al año; el ángulo $\beta$ con el norte cumple $\tan\beta =
16.6/41.9 \approx 0.40$, luego $\beta \approx \ang{22}$: la India avanza hacia
Eurasia en dirección norte desviada \ang{22} hacia el este. (2) \qty{45}{km}
por millón de años, luego \qty{450}{km} en 10 millones de años.
\end{solution}

\begin{solution}{exo:g11:plate-motions:10}
$(2432 - 1058)/(27.7 - 10.3) = 1374/17.4 \approx \qty{79}{km}$ por millón
de años, unos \qty{7.9}{cm} al año, algo menos que \qty{9.1}{cm} al año.
Cada edad tiene una incertidumbre de alrededor de medio millón de años; la
diferencia de edades, $17.4 \pm 1.0$ millones de años, da una velocidad de
entre unos 75 y \qty{84}{km} por millón de años. Entre estos dos volcanes la
placa parece haberse movido algo más despacio que de media.
\end{solution}

\begin{solution}{exo:g11:plate-motions:11}
(1) $\vec v_{A/C} = \vec v_{A/B} + \vec v_{B/C}$: \qty{3}{cm} al año hacia el
este más \qty{4}{cm} al año hacia el norte. Las dos flechas forman un ángulo
recto, y el tercer lado del triángulo mide $\sqrt{3^2 + 4^2} = \qty{5}{cm}$ al
año, en dirección norte desviada \ang{37} hacia el este ($\tan\beta = 3/4$).
(2) $\vec v_{C/A}
= -\vec v_{A/C}$: \qty{5}{cm} al año en la dirección opuesta, sur desviada
\ang{37} hacia el oeste.
\end{solution}

\begin{solution}{exo:g11:plate-motions:12}
(1) $54/0.375 = \qty{144}{km}$ por millón de años, \qty{14.4}{cm} al año. (2)
De $29/0.375 \approx 77$ a $79/0.375 \approx \qty{211}{km}$ por millón de
años: entre unos 8 y \qty{21}{cm} al año. (3) Cerca del punto caliente, la
distancia no es mayor que el tamaño de un volcán o de la región en la que el
punto caliente alimenta volcanes, y la edad no es mayor que el tiempo que
tarda en crecer un volcán: ambas son muy inciertas en términos relativos. Lejos
del punto caliente, un error de \qty{25}{km} sobre \qty{2400}{km}, o de medio
millón de años sobre 28 millones, es solo del uno o el dos por ciento.
\end{solution}

\begin{solution}{exo:g11:plate-motions:13}
(1) Diez veces la dispersión son \qty{50}{mm}: $50/72 \approx 0.7$ años, unos
ocho meses. (2) $50/5 = 10$ años. (3) A lo largo de millones de años las
placas se desplazan de decenas a cientos de kilómetros, distancias que
registran las bandas magnéticas y que eran fáciles de medir en el mar; un
posicionamiento con una precisión de pocos milímetros, que hace visible un
movimiento en meses, llegó mucho más tarde.
\end{solution}

\begin{solution}{exo:g11:plate-motions:14}
(1) En milímetros al año,
\[
  \vec v_{\mathrm{SA/AF}} = (-4.9 - 27.1,\ 11.0 - 17.9) = (-32.0, -6.9).
\] Su módulo es $\sqrt{32.0^2 + 6.9^2} \approx
\qty{32.7}{mm}$ al año, y $\tan\alpha = 6.9/32.0 \approx 0.22$ da una
dirección oeste desviada \ang{12} hacia el sur. (2) La \omterm{def:g11:drift-to-spreading:spreading-rate}{velocidad de expansión
total} es de unos \qty{3.3}{cm} al año, y la semivelocidad, de unos
\qty{1.6}{cm} al año. (3) La semivelocidad de los últimos millones de años,
\qty{1.6}{cm} al año, es algo menor que los \qty{1.9}{cm} al año promediados
sobre 66 millones de años: puede que el Atlántico Sur se haya abierto
recientemente algo más despacio que de media, o que la diferencia se deba a
las incertidumbres de las edades de los sedimentos. Las dos medidas coinciden
en que la dorsal es lenta.
\end{solution}

\begin{solution}{exo:g11:plate-motions:15}
(1) Una recta que sube \qty{20}{mm} al año, una caída brusca de
\qty{50}{mm} a los cinco años, y luego otra recta que vuelve a subir con la
misma pendiente. (2) La estación termina a $10 \times 20 - 50 = \qty{150}{mm}$
al este de su punto de partida: $150/10 = \qty{15}{mm}$ al año, una cuarta
parte menos de lo debido. (3) Cada mitad da \qty{20}{mm} al año. (4) El salto
es un movimiento aparte, brusco, causado por el terremoto: el geodesta mide su
tamaño con las posiciones de justo antes y justo después, lo resta de la serie
(o ajusta las dos partes con la misma pendiente y un escalón entre ellas), y
conserva \qty{20}{mm} al año como velocidad de la estación.
\end{solution}

\begin{solution}{pb:g11:plate-motions:1}
\textbf{1.} El Mauna Kea (\num{0.375} millones de años) es el más joven, y
Suiko (centro) (\num{64.7} millones de años), el más antiguo. Las edades
aumentan con la distancia a Kilauea.

\textbf{2.} $\qty{1}{km} = 10^6~\unit{mm}$ y un millón de años son $10^6$
años: \qty{1}{km} por millón de años es $10^6~\unit{mm}/10^6$ años, es decir,
\qty{1}{mm} al año.

\textbf{3.} $519/5.1 \approx \qty{102}{mm}$ al año.

\textbf{4.} Necker $1058/10.3 \approx 102.7$; Laysan $1818/19.9 \approx
91.4$; Midway $2432/27.7 \approx 87.8$; Colahan $3128/38.6 \approx 81.0$;
Abbott $3280/38.7 \approx 84.8$ (milímetros al año).

\textbf{5.} Los puntos (54; 0.375), (519; 5.1), (780; 7.2), (1058; 10.3),
(1818; 19.9), (2432; 27.7), (3128; 38.6), (3280; 38.7) quedan cerca de una
recta que pasa por el origen, con cierta dispersión (respecto de una recta de
\qty{91}{km} por millón de años, Kauai, Nihoa y Necker quedan un poco por
debajo, Laysan sobre ella, y Colahan y Abbott un poco por encima).

\textbf{6.} Cada volcán se construyó sobre el punto caliente y desde entonces
se ha movido con la placa. En el tiempo $t$, la placa se ha desplazado $v\,t$
en una dirección constante, así que el volcán está a $d = v\,t$ del punto que
ahora está sobre el punto caliente: $t = d/v$.

\textbf{7.} $(102.7 + 91.4 + 87.8 + 81.0 + 84.8)/5 \approx \qty{89.5}{mm}$ al
año.

\textbf{8.} 89.5 frente a \qty{91}{km} por millón de años: difieren en menos
del \qty{2}{\percent}.

\textbf{9.} $519/91 \approx 5.7$ millones de años, frente a una edad de
\num{5.1} millones de años: valores cercanos; Kauai es algo más joven de lo que
predice la recta media.

\textbf{10.} Los volcanes son cada vez más antiguos hacia el oeste-noroeste:
la placa los transporta en esa dirección. La placa del Pacífico se mueve hacia
el oeste-noroeste sobre el punto caliente.

\textbf{11.} $(4860 - 3493)/(64.7 - 42.4) = 1367/22.3 \approx \qty{61}{km}$
por millón de años, unos \qty{6.1}{cm} al año.

\textbf{12.} $65/91 \approx 0.71$: a lo largo de la cadena del Emperador la
placa se movió alrededor de un \qty{30}{\percent} más despacio que a lo largo
de la cadena de Hawái.

\textbf{13.} $3668/50 \approx \qty{73}{km}$ por millón de años, unos
\qty{7.3}{cm} al año, menos que \qty{91}{km} por millón de años. La recta por el origen
usa las edades de 1987, que sitúan el final de la cadena de Hawái entre 39 y
43 millones de años; la edad más reciente, de unos 50 millones de años, para
el comienzo del codo reparte la misma distancia en un tiempo más largo. Las
dos son correctas si la placa se movió más despacio en la primera parte de la
cadena de Hawái y más deprisa después: la velocidad no fue constante.

\textbf{14.} Un cambio en la dirección del movimiento de la placa del
Pacífico, hace unos 50 millones de años.

\textbf{15.} La dirección de la cadena es la dirección en la que se movió la
placa; el espaciado de las edades a lo largo de ella es su velocidad. Una
placa puede acelerar o frenar sin girar, o girar sin cambiar de velocidad, así
que las dos cosas son independientes. Los cambios de velocidad sin cambio de
dirección a lo largo de la cadena del Emperador muestran que su registro de
velocidades y el codo son sucesos distintos: el codo no se debió a que el
punto caliente se frenara.

\textbf{16.} $\sqrt{63.2^2 + 34.5^2} \approx \qty{72.0}{mm}$ al año, y
$\tan\alpha = 34.5/63.2 \approx 0.55$ da una dirección oeste desviada \ang{29}
hacia el norte.

\textbf{17.} El oeste-noroeste está a \ang{22.5} al norte del oeste, y el
noroeste a \ang{45}: la dirección de la estación, \ang{29} al norte del oeste,
queda entre las dos, cerca de la dirección de la cadena. Es coherente.

\textbf{18.} La velocidad de la estación se mide respecto de un marco en el
que la \omterm{def:g11:lithosphere-asthenosphere:lithosphere}{litosfera} en su conjunto no gira, mientras que la cadena da el
movimiento respecto del punto caliente: si los dos marcos se mueven uno
respecto del otro, las velocidades difieren. Además, la estación mide los
últimos 30 años, y la cadena los últimos 40 millones de años, durante los
cuales la velocidad puede haber cambiado. La dispersión de las edades deja
también una incertidumbre de unos pocos milímetros al año en la velocidad de
la cadena.

\textbf{19.} $100/91 \approx 1.1$ millones de años.

\textbf{20.} \textbf{La placa del Pacífico se ha movido sobre el punto caliente
de Hawái a unos \qty{9}{cm} al año} durante los últimos 40 millones de años,
hacia el oeste-noroeste; antes del codo se movía más despacio, y en otra
dirección.
\end{solution}
